\documentclass[10pt,a4paper]{amsart}
\usepackage[margin=19mm]{geometry}
\usepackage{amsmath,amssymb,mathtools}
\usepackage[scr=rsfs,cal=euler]{mathalfa}
\usepackage{xfrac}
\usepackage[unicode,hidelinks,bookmarks=false]{hyperref}
\newcommand{\ie}{i.e.\ }
\newcommand{\cf}{cf.\ }
\newcommand{\resp}{respectively\ }
\newcommand{\Subsection}{\subsection}
\providecommand{\nobreakdash}{\nobreak\hbox{-}}
\setlength{\emergencystretch}{2em}
\multlinegap0pt
\theoremstyle{plain}
\newtheorem{theorem}{Theorem}
\newtheorem{proposition}{Proposition}
\newtheorem{lemma}{Lemma}
\newtheorem{corollary}{Corollary}
\newtheorem{condition}{Condition}
\newtheorem{definition}{Definition}
\theoremstyle{remark}
\newtheorem{remark}{Remark}
\newcommand{\btheta}{\boldsymbol \theta}
\newcommand{\bTheta}{\boldsymbol \Theta}
\newcommand{\btc}{{\rm c}}
\def\T{{\scriptscriptstyle \top} }
\title[Autoregressive networks with dependent edges]{Autoregressive networks with dependent edges: consistency of the initial estimator for the global and local parameters}
\author{Jinyuan Chang, Qin Fang, Eric D. Kolaczyk, Peter W. MacDonald, Qiwei Yao}
\date{Source: arXiv:2404.15654v6 (CC BY 4.0)}
\begin{document}
\maketitle
\noindent\emph{Source and licence.} Autoregressive networks with dependent edges. Version arXiv:2404.15654v6 (CC BY 4.0), distributed under the Creative Commons Attribution 4.0 International licence, \url{http://creativecommons.org/licenses/by/4.0/}, as recorded in the arXiv licence metadata for this exact version and shown on the abstract page \url{https://arxiv.org/abs/2404.15654v6}. Full licence notice: licenses/arxiv-2404.15654-CC-BY-4.0.txt.\par\smallskip

\noindent\textit{Editorial scope.} Counted unit: the article's Theorem with source label \texttt{tm:consistency} (source0.tex lines 945--947), the consistency of the initial estimator for the global and for the local parameters of the autoregressive network model with dependent edges, delivered together with the article's complete original proof of it (source0.tex lines 3163--3227, the article's own proof subsection of 65 lines, whose heading names the counted theorem and whose argument closes with the article's completion sentence and end-of-proof symbol). No mathematics is written, repaired or re-derived: the counted statement and its proof are the article's own text line for line, in the article's own notation (the conditions on the network model, the sets of indices, the two rates, the uniform-covariance lemma that the proof states and then uses), and no retained line is substituted, re-flowed or omitted. Required context retained uncounted, in source order: the three conditions the counted statement invokes through the identification proposition (lines 687--699, 733--738 and 744--754); the statement of the article's identification proposition (lines 781--789), whose conclusion the counted statement invokes; the heading of the article's technical-proofs section (line 3041) and the headings of the appendix subsection and subsubsection in which the auxiliary lemma's own proof is given (lines 3480--3482), which the counted proof cites by section reference. The auxiliary lemma's own proof and the proofs of the identification proposition and of the three conditions are not delivered. Nothing else of the article is retained: its other theorems, its remaining lemmas, its data analysis, its supplementary material and all of its figures are omitted. No citation command occurs in any retained line, so the article's bibliography is not delivered and the wrapper carries no bibliography entry. Every cross-reference inside the retained spans resolves inside the delivered document. Editorial formatting only: an AMS \texttt{amsart} preamble that restates the environments and macros the retained lines use, namely the article's declarations of Theorem, Proposition, Lemma, Corollary, Condition, Definition and Remark, and the article's four macros for bold Greek parameters, for the complement of an index set and for the transpose symbol. The article is typeset in its own \texttt{article}-class format with its own fonts; no class or style file is shipped with this package and the wrapper is an amsart document. Numbering differs from the article, because the article numbers its environments by type without a section reset and because only a subset of environments is retained: the retained environments print with the wrapper's continuous numbering in their source order, so the three retained conditions print as Condition 1 to Condition 3, the retained identification proposition as Proposition 1, and the counted theorem as Theorem 1; the retained appendix headings print with the wrapper's section numbering, so the counted proof prints under Section 1 and the cited auxiliary-proof section prints as Section 1.2.1. The article's cross-references are not rewritten. No picture environment and no graphical inclusion occurs in any retained line, and the package delivers no image asset. One bundled source file, \texttt{sources/158085/source0.tex}, accompanies the delivered item as the exact source of every retained span. Every retained span is recorded in \texttt{delivery-evidence.json} with method \texttt{exact\_tex}.
\begin{condition}\label{as:gambd}
{\rm(i)} There exists some universal constant $C_{1}>0$ such that
\[
\min_{t\in[n]\setminus[m]}\min_{i,j:\,1\leq i<j\leq p}\inf_{\btheta\in\bTheta}\gamma_{i,j}^{t-1}(\btheta)\{1-\gamma_{i,j}^{t-1}(\btheta)\}\geq C_1\,.
\]
{\rm (ii)} For any $1\leq i<j\leq p$ and $t\in[n]\setminus[m]$, $\gamma_{i,j}^{t-1}(\btheta)$ is thrice continuously
differentiable with respect to $\btheta\in\bTheta$. Furthermore, there exists some universal constant
$C_{2}>0$ such that
\[
\max_{t\in[n]\setminus[m]}\max_{i,j:\,1\leq i<j\leq p}\sup_{\btheta\in\bTheta}\bigg|\frac{\partial^k \gamma_{i,j}^{t-1}(\btheta)}{\partial\btheta^k}\bigg|_\infty\leq C_2
\]
for any $k\in[3]$.
\end{condition}

\begin{condition}\label{as:bdSij}
There exists a universal constant $s\geq1$ such that
$
\max_{1\leq i<j\leq p}|\mathcal{I}_{i,j}|\leq s
$.
\end{condition}

\begin{condition}\label{as:eigen}
There exists a universal constant $C_3>0$ such that
\[
\min_{i,j:\,1\leq i<j\leq p}\lambda_{\min}\bigg\{\frac{1}{n-m}\sum_{t=m+1}^n\frac{\partial\gamma_{i,j}^{t-1}(\btheta_0)}{\partial\btheta_{\mathcal{I}_{i,j}}}\frac{\partial\gamma_{i,j}^{t-1}(\btheta_0)}{\partial\btheta_{\mathcal{I}_{i,j}}^{\T}}\bigg\}\geq C_3
\]
%and
%\[
%\max_{i,j:\,i\neq j}\lambda_{\max}\bigg\{\frac{1}{n-m}\sum_{t=m+1}^n\frac{\partial\gamma_{i,j}^{t-1}(\btheta_0)}{\partial\btheta_{\mathcal{I}_{i,j}}}\frac{\partial\gamma_{i,j}^{t-1}(\btheta_0)}{\partial\btheta_{\mathcal{I}_{i,j}}^{\T}}\bigg\}\leq C_4
%\]
with probability approaching one when $n\to \infty$. 
\end{condition}

\begin{proposition}\label{pn:ident}
Let Conditions {\rm\ref{as:gambd}--\ref{as:eigen}} hold, and $C_*=2(2C_1^{-2}+C_1^{-3})C_2^3+3(C_1^{-1}+C_1^{-2})C_2^2+C_1^{-1}C_2$ with $(C_1,C_2)$ specified in Condition {\rm\ref{as:gambd}}. Assume $\sup_{\btheta\in\bTheta}|\btheta-\btheta_0|_\infty<2C_3/(C_*s^3)$.
As $n\rightarrow\infty$, it holds 
with probability approaching one that 
\begin{align*}
\ell_{n,p}^{(l)}(\btheta_0)-\ell_{n,p}^{(l)}(\btheta)\geq \frac{\bar{C}}{|\mathcal{S}_l|}\sum_{(i,j)\in\mathcal{S}_l}|\btheta_{\mathcal{I}_{i,j}}-\btheta_{0,\mathcal{I}_{i,j}}|_2^2
\end{align*}
for any $\btheta\in\bTheta$ and $l\in[q]$, where $\bar{C}>0$ is a universal constant.  
\end{proposition}

\begin{theorem}\label{tm:consistency}
Let the conditions of Proposition {\rm\ref{pn:ident}} hold. Then $|\tilde{\btheta}_{\mathcal{G}}-\btheta_{0,\mathcal{G}}|_2=O_{\rm p}(c_{n,\mathcal{G}})$ and $|\tilde{\btheta}_{\mathcal{G}^{\btc}}-\btheta_{0,\mathcal{G}^{\btc}}|_\infty=O_{\rm p}(c_{n,\mathcal{G}^{\btc}})$. 
\end{theorem}

\section{Technical proofs} \label{app:proofs}

\subsection{Proof of Theorem \ref{tm:consistency}} \label{pr:tm:consistency}
Recall $S_{\mathcal{H},\min}=\min_{l\in\mathcal{H}}|\mathcal{S}_l|$ for any $\mathcal{H}\subset[q]$, and \begin{align*}
\left\{\begin{aligned}
   c_{n,\mathcal{G}}^2=&~\frac{q\log(n|\mathcal{S}_{l'}|)}{\sqrt{n|\mathcal{S}_{l'}|}}+\frac{q^{3/2}\log^{3/2}(n|\mathcal{S}_{l'}|)}{\sqrt{n}|\mathcal{S}_{l'}|}\,, \\
   c_{n,\mathcal{G}^{\btc}}^2=&~\frac{q\log(nS_{\mathcal{G}^{\btc},\min})}{\sqrt{nS_{\mathcal{G}^{\btc},\min}}}+\frac{q^{3/2}\log^{3/2}(nS_{\mathcal{G}^{\btc},\min})}{\sqrt{n}S_{\mathcal{G}^{\btc},\min}}\,,
\end{aligned}\right.
\end{align*}
where $l'\in\mathcal{G}$. 
To show Theorem \ref{tm:consistency}, we first present the following lemma whose proof is given in Section \ref{sec:pflauniformcov}. 

\begin{lemma}\label{la:uniformcov}
 Assume Conditions {\rm\ref{as:gambd}} and {\rm\ref{as:bdSij}} hold. Then  
\begin{align*}
\max_{l\in\mathcal{H}}\sup_{\btheta \in \bTheta}|\hat{\ell}_{n,p}^{(l)}(\btheta) - \ell_{n,p}^{(l)}(\btheta)|=O_{\rm p}\bigg\{\frac{q\log(nS_{\mathcal{H},\min})}{\sqrt{nS_{\mathcal{H},\min}}}\bigg\}+O_{\rm p}\bigg\{\frac{q^{3/2}\log^{3/2}(nS_{\mathcal{H},\min})}{\sqrt{n}S_{\mathcal{H},\min}}\bigg\}
\end{align*}
for any $\mathcal{H}\subset[q]$.
\end{lemma}

%\begin{lemma}\label{la:inirate}
%Assume Conditions {\rm\ref{as:gambd}}--{\rm\ref{as:eigen}} hold. Then $|\widetilde{\btheta}_{\mathcal{G}}-\btheta_{0,\mathcal{G}}|_2=O_{\rm p}(c_{n,\mathcal{G}})$ and $|\widetilde{\btheta}_{\mathcal{G}^{\btc}}-\btheta_{0,\mathcal{G}^{\btc}}|_\infty=O_{\rm p}(c_{n,\mathcal{G}^{\btc}})$. 
%\end{lemma}


Notice that $
\hat{\btheta}_*^{(l)}=(\hat \theta_{*,1}^{(l)}, \dots,\hat \theta_{*,q}^{(l)} )^{\T}=\arg\max_{\btheta\in\bTheta}\hat{\ell}_{n,p}^{(l)}(\btheta)$ for any $l\in[q]$. Then
\begin{align*}
{\ell}_{n,p}^{(l')}(\btheta_0)-\sup_{\btheta\in\bTheta}|\hat{\ell}_{n,p}^{(l')}(\btheta)-{\ell}_{n,p}^{(l')}(\btheta)|\leq&~\hat{\ell}_{n,p}^{(l')}(\btheta_0)\\
\leq&~\hat{\ell}_{n,p}^{(l')}(\hat{\btheta}_*^{(l')})\leq {\ell}_{n,p}^{(l')}(\hat{\btheta}_*^{(l')})+\sup_{\btheta\in\bTheta}|\hat{\ell}_{n,p}^{(l')}(\btheta)-{\ell}_{n,p}^{(l')}(\btheta)|\,,
\end{align*}
which implies
\begin{align*}%\label{eq:diffbd1}
0\leq {\ell}_{n,p}^{(l')}(\btheta_0)-{\ell}_{n,p}^{(l')}(\hat{\btheta}_*^{(l')})\leq 2\sup_{\btheta\in\bTheta}|\hat{\ell}_{n,p}^{(l')}(\btheta)-{\ell}_{n,p}^{(l')}(\btheta)|\,.
\end{align*}
Recall $\tilde{\btheta}_{\mathcal{G}}=\hat{\btheta}_{*,\mathcal{G}}^{(l')}$. Selecting $\mathcal{H}=\{l'\}$ in Lemma \ref{la:uniformcov}, we have $\sup_{\btheta\in\bTheta}|\hat{\ell}_{n,p}^{(l')}(\btheta)-{\ell}_{n,p}^{(l')}(\btheta)|=O_{\rm p}(c_{n,\mathcal{G}}^2)$, which implies 
\begin{align}\label{eq:diffbd11}
{\ell}_{n,p}^{(l')}(\btheta_0)-{\ell}_{n,p}^{(l')}(\hat{\btheta}_*^{(l')})=O_{\rm p}(c_{n,\mathcal{G}}^2)\,.
\end{align}
For any diverging $\epsilon_{n,p}>0$, if $|\tilde{\btheta}_{\mathcal{G}}-\btheta_{0,\mathcal{G}}|_2\geq \epsilon_{n,p}c_{n,\mathcal{G}}$, Proposition \ref{pn:ident} yields that
\[
{\ell}_{n,p}^{(l')}(\btheta_0)-{\ell}_{n,p}^{(l')}(\hat{\btheta}_*^{(l)})\geq\bar{C}\epsilon_{n,p}^2c_{n,\mathcal{G}}^2 
\]
with probability approaching one, which contradicts with \eqref{eq:diffbd11} and then implies $|\tilde{\btheta}_{\mathcal{G}}-\btheta_{0,\mathcal{G}}|_2=O_{\rm p}( \epsilon_{n,p}c_{n,\mathcal{G}})$. Notice that we can select arbitrary slowly diverging $\epsilon_{n,p}$. Following a standard result from
probability theory, we have $|\tilde{\btheta}_{\mathcal{G}}-\btheta_{0,\mathcal{G}}|_2=O_{\rm p}(c_{n,\mathcal{G}})$. 

For any $l\in\mathcal{G}^{\btc}$, %since $
%\widehat{\btheta}_*^{(l)}=(\hat \theta_{*,1}^{(l)}, \dots,\hat \theta_{*,q}^{(l)} )^{\T}=\arg\max_{\btheta\in\bTheta}\hat{\ell}_{n,p}^{(l)}(\btheta)$, then
%\begin{align*}
%{\ell}_{n,p}^{(l)}(\btheta_0)-\sup_{\btheta\in\bTheta}|\hat{\ell}_{n,p}^{(l)}(\btheta)-{\ell}_{n,p}^{(l)}(\btheta)|\leq&~\hat{\ell}_{n,p}^{(l)}(\btheta_0)\\
%\leq&~\hat{\ell}_{n,p}^{(l)}(\widehat{\btheta}_*^{(l)})\leq {\ell}_{n,p}^{(l)}(\widehat{\btheta}_*^{(l)})+\sup_{\btheta\in\bTheta}|\hat{\ell}_{n,p}^{(l)}(\btheta)-{\ell}_{n,p}^{(l)}(\btheta)|\,,
%\end{align*}
%which implies
we also have 
\begin{align*}%\label{eq:diffbd1}
0\leq {\ell}_{n,p}^{(l)}(\btheta_0)-{\ell}_{n,p}^{(l)}(\hat{\btheta}_*^{(l)})\leq 2\sup_{\btheta\in\bTheta}|\hat{\ell}_{n,p}^{(l)}(\btheta)-{\ell}_{n,p}^{(l)}(\btheta)|\,.
\end{align*}
Recall $\tilde{\btheta}_{\mathcal{G}^{\btc}}=(\hat{\theta}_{*,l}^{(l)})_{l\in\mathcal{G}^{\btc}}$. Selecting $\mathcal{H}=\mathcal{G}^{\btc}$ in Lemma \ref{la:uniformcov}, we have $\max_{l\in\mathcal{G}^{\btc}}\sup_{\btheta\in\bTheta}|\hat{\ell}_{n,p}^{(l)}(\btheta)-{\ell}_{n,p}^{(l)}(\btheta)|=O_{\rm p}(c_{n,\mathcal{G}^{\btc}}^2)$, which implies 
\begin{align}\label{eq:diffbd1}
\max_{l\in\mathcal{G}^{\btc}}\big\{{\ell}_{n,p}^{(l)}(\btheta_0)-{\ell}_{n,p}^{(l)}(\hat{\btheta}_*^{(l)})\big\}=O_{\rm p}(c_{n,\mathcal{G}^{\btc}}^2)\,.
\end{align}
For any diverging $\epsilon_{n,p}>0$, if $|\tilde{\btheta}_{\mathcal{G}^{\btc}}-\btheta_{0,\mathcal{G}^{\btc}}|_\infty\geq \epsilon_{n,p}c_{n,\mathcal{G}^{\btc}}$, Proposition \ref{pn:ident} yields that
\[
\max_{l\in\mathcal{G}^{\btc}}\big\{{\ell}_{n,p}^{(l)}(\btheta_0)-{\ell}_{n,p}^{(l)}(\hat{\btheta}_*^{(l)})\big\}\geq\bar{C}\epsilon_{n,p}^2c_{n,\mathcal{G}^{\btc}}^2 
\]
with probability approaching one, which contradicts with \eqref{eq:diffbd1} and then implies $|\tilde{\btheta}_{\mathcal{G}^{\btc}}-\btheta_{0,\mathcal{G}^{\btc}}|_\infty=O_{\rm p}( \epsilon_{n,p}c_{n,\mathcal{G}^{\btc}})$. Notice that we can select arbitrary slowly diverging $\epsilon_{n,p}$. Following a standard result from
probability theory, we have $|\tilde{\btheta}_{\mathcal{G}^{\btc}}-\btheta_{0,\mathcal{G}^{\btc}}|_\infty=O_{\rm p}(c_{n,\mathcal{G}^{\btc}})$. We complete the proof of Theorem \ref{tm:consistency}.  $\hfill\Box$

\subsection{Proofs of auxiliary lemmas}

\subsubsection{Proof of Lemma \ref{la:uniformcov}}\label{sec:pflauniformcov}

\end{document}
